\textbf{Design ablation (H5 and normalisation; 400 steps).} Table~\ref{tab:abl} removes the decomposition (DLinear only) or the instance normalisation (each learned model) at $H=96$; ``Full'' is the main-grid row. \emph{H5 is supported}: a single linear map without decomposition leaves DLinear's MSE almost unchanged: \rhval{abl_design/linear-no-decomp/trend/mse_96/mean:3} against \rhval{main/dlinear/trend/mse_96/mean:3} (trend), \rhval{abl_design/linear-no-decomp/regime/mse_96/mean:3} against \rhval{main/dlinear/regime/mse_96/mean:3} (regime) and \rhval{abl_design/linear-no-decomp/etth1/mse_96/mean:3} against \rhval{main/dlinear/etth1/mse_96/mean:3} (etth1), each change within 5\%. Our trends are mild and the window normalisation already removes the level, so this does not show the decomposition is useless in general.

Instance normalisation matters more, and not in one direction. On the trend task removing it raises the seed-mean MSE of DLinear from \rhval{main/dlinear/trend/mse_96/mean:3} to \rhval{abl_design/dlinear-no-norm/trend/mse_96/mean:3}, of PatchTST from \rhval{main/patchtst/trend/mse_96/mean:3} to \rhval{abl_design/patchtst-no-norm/trend/mse_96/mean:3} and of the GRU from \rhval{main/gru/trend/mse_96/mean:3} to \rhval{abl_design/gru-no-norm/trend/mse_96/mean:3}, but all three means are driven by one dataset (seed 1, whose trend MSE is several times that of the other two seeds). The median over seeds moves from \rhval{main/dlinear/trend/mse_96/median:4} to \rhval{abl_design/dlinear-no-norm/trend/mse_96/median:4} for DLinear, from \rhval{main/patchtst/trend/mse_96/median:4} to \rhval{abl_design/patchtst-no-norm/trend/mse_96/median:4} for PatchTST and from \rhval{main/gru/trend/mse_96/median:4} to \rhval{abl_design/gru-no-norm/trend/mse_96/median:4} for the GRU: the two nonlinear models degrade strongly on a typical seed, DLinear only mildly. On regime, in contrast, removing the normalisation \emph{lowers} the seed-mean MSE for all three models (\rhval{main/dlinear/regime/mse_96/mean:3} to \rhval{abl_design/dlinear-no-norm/regime/mse_96/mean:3} DLinear, \rhval{main/patchtst/regime/mse_96/mean:3} to \rhval{abl_design/patchtst-no-norm/regime/mse_96/mean:3} PatchTST, \rhval{main/gru/regime/mse_96/mean:3} to \rhval{abl_design/gru-no-norm/regime/mse_96/mean:3} GRU; the seed medians in Table~\ref{tab:abl} also fall), and on etth1 it hurts the GRU (\rhval{main/gru/etth1/mse_96/mean:3} to \rhval{abl_design/gru-no-norm/etth1/mse_96/mean:3}) and changes the other two little (\rhval{main/dlinear/etth1/mse_96/mean:3} to \rhval{abl_design/dlinear-no-norm/etth1/mse_96/mean:3} and \rhval{main/patchtst/etth1/mse_96/mean:3} to \rhval{abl_design/patchtst-no-norm/etth1/mse_96/mean:3}). The regime result is exploratory (not a registered hypothesis); window-statistics normalisation may hurt when the signal scale changes between regimes, but we did not test that.

\begin{table}[t]\centering\caption{Design ablation at $H=96$, 400 steps (test MSE). ``Full'' is the main-grid model and ``No norm'' the same model without instance normalisation, each as the mean and as the median over the three seeds; the last column is DLinear without decomposition (seed mean).}\label{tab:abl}
\resizebox{\columnwidth}{!}{\begin{tabular}{llrrrrr}
\toprule
 & & \multicolumn{2}{c}{seed mean} & \multicolumn{2}{c}{seed median} & No decomp\\
Task & System & Full & No norm & Full & No norm & (seed mean)\\
\midrule
trend & DLinear & \rhval{main/dlinear/trend/mse_96/mean:3} & \rhval{abl_design/dlinear-no-norm/trend/mse_96/mean:3} & \rhval{main/dlinear/trend/mse_96/median:4} & \rhval{abl_design/dlinear-no-norm/trend/mse_96/median:4} & \rhval{abl_design/linear-no-decomp/trend/mse_96/mean:3}\\
trend & PatchTST & \rhval{main/patchtst/trend/mse_96/mean:3} & \rhval{abl_design/patchtst-no-norm/trend/mse_96/mean:3} & \rhval{main/patchtst/trend/mse_96/median:4} & \rhval{abl_design/patchtst-no-norm/trend/mse_96/median:4} & --\\
trend & GRU & \rhval{main/gru/trend/mse_96/mean:3} & \rhval{abl_design/gru-no-norm/trend/mse_96/mean:3} & \rhval{main/gru/trend/mse_96/median:4} & \rhval{abl_design/gru-no-norm/trend/mse_96/median:4} & --\\
regime & DLinear & \rhval{main/dlinear/regime/mse_96/mean:3} & \rhval{abl_design/dlinear-no-norm/regime/mse_96/mean:3} & \rhval{main/dlinear/regime/mse_96/median:4} & \rhval{abl_design/dlinear-no-norm/regime/mse_96/median:4} & \rhval{abl_design/linear-no-decomp/regime/mse_96/mean:3}\\
regime & PatchTST & \rhval{main/patchtst/regime/mse_96/mean:3} & \rhval{abl_design/patchtst-no-norm/regime/mse_96/mean:3} & \rhval{main/patchtst/regime/mse_96/median:4} & \rhval{abl_design/patchtst-no-norm/regime/mse_96/median:4} & --\\
regime & GRU & \rhval{main/gru/regime/mse_96/mean:3} & \rhval{abl_design/gru-no-norm/regime/mse_96/mean:3} & \rhval{main/gru/regime/mse_96/median:4} & \rhval{abl_design/gru-no-norm/regime/mse_96/median:4} & --\\
etth1 & DLinear & \rhval{main/dlinear/etth1/mse_96/mean:3} & \rhval{abl_design/dlinear-no-norm/etth1/mse_96/mean:3} & \rhval{main/dlinear/etth1/mse_96/median:4} & \rhval{abl_design/dlinear-no-norm/etth1/mse_96/median:4} & \rhval{abl_design/linear-no-decomp/etth1/mse_96/mean:3}\\
etth1 & PatchTST & \rhval{main/patchtst/etth1/mse_96/mean:3} & \rhval{abl_design/patchtst-no-norm/etth1/mse_96/mean:3} & \rhval{main/patchtst/etth1/mse_96/median:4} & \rhval{abl_design/patchtst-no-norm/etth1/mse_96/median:4} & --\\
etth1 & GRU & \rhval{main/gru/etth1/mse_96/mean:3} & \rhval{abl_design/gru-no-norm/etth1/mse_96/mean:3} & \rhval{main/gru/etth1/mse_96/median:4} & \rhval{abl_design/gru-no-norm/etth1/mse_96/median:4} & --\\
\bottomrule
\end{tabular}}\end{table}

\textbf{Regime-length and noise sweeps.} Tables~\ref{tab:dwell} and \ref{tab:noise} and Fig.~\ref{fig:sweeps} (regime task, $H=96$, 3 seeds per cell) are discussed under H2 and H3. Three cautions apply. At dwell 4000 a series of 8000 steps contains at most a few regime changes, so that cell depends on which regimes fall in the test split, and its spread across seeds is large (Fig.~\ref{fig:sweeps}). The single 400-step sweep win for PatchTST at dwell 1000 (\rhval{sweep_dwell/patchtst@dwell=1000/regime/mse_96/mean:3} against \rhval{sweep_dwell/dlinear@dwell=1000/regime/mse_96/mean:3} for DLinear) is not robust to the budget: at 2000 steps the same cell is \rhval{sweep_dwell_2000/patchtst@dwell=1000/regime/mse_96/mean:3} against \rhval{sweep_dwell_2000/dlinear@dwell=1000/regime/mse_96/mean:3}, better in only 2/3 seeds. The dwell and noise axes were each swept with the other at its default, so interactions are unknown.

\begin{table}[t]\centering\caption{Dwell sweep (regime task, $H=96$) at 400 steps (top) and 2000 steps (bottom): seed-mean MSE; superscript = number of seeds (of 3) in which the nonlinear model is better than DLinear; $^\ast$ = win. Dwell 500 is the main-grid row (top) or the budget-sweep row (bottom).}\label{tab:dwell}
\begin{tabular}{lrrrrrrr}
\toprule
dwell & DLinear & \multicolumn{3}{c}{PatchTST} & \multicolumn{3}{c}{GRU}\\
 & MSE & MSE & vs DL (\%) & seeds & MSE & vs DL (\%) & seeds\\
\midrule
100 & 0.774 & 0.773 & -0.2 & 1/3 & 0.743 & -4.1 & 3/3\\
250 & 0.554 & 0.529 & -4.5 & 2/3 & 0.537 & -3.2 & 2/3\\
500 & 0.434 & 0.413 & -4.8 & 3/3 & 0.446 & +2.9 & 1/3\\
1000 & 0.272 & 0.258 & -5.1$^\ast$ & 3/3 & 0.337 & +23.8 & 0/3\\
4000 & 0.312 & 0.307 & -1.7 & 3/3 & 0.599 & +91.7 & 0/3\\
\bottomrule
\end{tabular}\\[3pt]
\begin{tabular}{lrrr}
\toprule
dwell & DLinear & PatchTST & GRU\\
\midrule
100 & \rhval{sweep_dwell_2000/dlinear@dwell=100/regime/mse_96/mean:3} & \rhval{sweep_dwell_2000/patchtst@dwell=100/regime/mse_96/mean:3}$^{0}$ & \rhval{sweep_dwell_2000/gru@dwell=100/regime/mse_96/mean:3}$^{3\ast}$\\
250 & \rhval{sweep_dwell_2000/dlinear@dwell=250/regime/mse_96/mean:3} & \rhval{sweep_dwell_2000/patchtst@dwell=250/regime/mse_96/mean:3}$^{2}$ & \rhval{sweep_dwell_2000/gru@dwell=250/regime/mse_96/mean:3}$^{2}$\\
500 & \rhval{sweep_steps/dlinear@steps=2000/regime/mse_96/mean:3} & \rhval{sweep_steps/patchtst@steps=2000/regime/mse_96/mean:3}$^{3}$ & \rhval{sweep_steps/gru@steps=2000/regime/mse_96/mean:3}$^{3\ast}$\\
1000 & \rhval{sweep_dwell_2000/dlinear@dwell=1000/regime/mse_96/mean:3} & \rhval{sweep_dwell_2000/patchtst@dwell=1000/regime/mse_96/mean:3}$^{2}$ & \rhval{sweep_dwell_2000/gru@dwell=1000/regime/mse_96/mean:3}$^{2}$\\
4000 & \rhval{sweep_dwell_2000/dlinear@dwell=4000/regime/mse_96/mean:3} & \rhval{sweep_dwell_2000/patchtst@dwell=4000/regime/mse_96/mean:3}$^{1}$ & \rhval{sweep_dwell_2000/gru@dwell=4000/regime/mse_96/mean:3}$^{1}$\\
\bottomrule
\end{tabular}\end{table}
\begin{table}[t]\centering\caption{Noise sweep (regime task, $H=96$, 400 steps), same layout. Noise 0.3 is the main-grid row.}\label{tab:noise}
\begin{tabular}{lrrr}
\toprule
noise & DLinear & PatchTST & GRU\\
\midrule
0.1 & \rhval{sweep_noise/dlinear@noise=0.1/regime/mse_96/mean:3} & \rhval{sweep_noise/patchtst@noise=0.1/regime/mse_96/mean:3}$^{3\ast}$ & \rhval{sweep_noise/gru@noise=0.1/regime/mse_96/mean:3}$^{1}$\\
0.3 & \rhval{main/dlinear/regime/mse_96/mean:3} & \rhval{main/patchtst/regime/mse_96/mean:3}$^{3}$ & \rhval{main/gru/regime/mse_96/mean:3}$^{1}$\\
0.6 & \rhval{sweep_noise/dlinear@noise=0.6/regime/mse_96/mean:3} & \rhval{sweep_noise/patchtst@noise=0.6/regime/mse_96/mean:3}$^{3}$ & \rhval{sweep_noise/gru@noise=0.6/regime/mse_96/mean:3}$^{1}$\\
1.5 & \rhval{sweep_noise/dlinear@noise=1.5/regime/mse_96/mean:3} & \rhval{sweep_noise/patchtst@noise=1.5/regime/mse_96/mean:3}$^{3}$ & \rhval{sweep_noise/gru@noise=1.5/regime/mse_96/mean:3}$^{2}$\\
3 & \rhval{sweep_noise/dlinear@noise=3.0/regime/mse_96/mean:3} & \rhval{sweep_noise/patchtst@noise=3.0/regime/mse_96/mean:3}$^{3}$ & \rhval{sweep_noise/gru@noise=3.0/regime/mse_96/mean:3}$^{3}$\\
\bottomrule
\end{tabular}\end{table}
\begin{figure}[t]\centering\includegraphics[width=\columnwidth]{sweeps.pdf}\caption{Test MSE at $H=96$ on the regime task versus mean regime dwell (left) and noise sd (right, 400 steps only). Solid: 400 steps, mean $\pm$ sd over 3 seeds; dashed: 2000 steps, mean.}\label{fig:sweeps}\end{figure}
